\apendice{Especificación de Requisitos}

\section{Introducción}

Una muestra de cómo podría ser una tabla de casos de uso:

% Caso de Uso 1 -> Consultar Experimentos.
\begin{table}[p]
	\centering
	\begin{tabularx}{\linewidth}{ p{0.21\columnwidth} p{0.71\columnwidth} }
		\toprule
		\textbf{CU-1}    & \textbf{Ejemplo de caso de uso}\\
		\toprule
		\textbf{Versión}              & 1.0    \\
		\textbf{Autor}                & Alumno \\
		\textbf{Requisitos asociados} & RF-xx, RF-xx \\
		\textbf{Descripción}          & La descripción del CU \\
		\textbf{Precondición}         & Precondiciones (podría haber más de una) \\
		\textbf{Acciones}             &
		\begin{enumerate}
			\def\labelenumi{\arabic{enumi}.}
			\tightlist
			\item Pasos del CU
			\item Pasos del CU (añadir tantos como sean necesarios)
		\end{enumerate}\\
		\textbf{Postcondición}        & Postcondiciones (podría haber más de una) \\
		\textbf{Excepciones}          & Excepciones \\
		\textbf{Importancia}          & Alta o Media o Baja... \\
		\bottomrule
	\end{tabularx}
	\caption{CU-1 Nombre del caso de uso.}
\end{table}

\section{Objetivos generales}

\section{Catálogo de requisitos}
\subsection{Requisitos Funcionales}
\begin{enumerate}
	\item RF-01. Autenticación de usuarios.
	\item RF-02. Consultar reservas.
	\item RF-03. Gestionar reservas.
	\begin{enumerate}
		\item RF-03.1. Crear reserva.
		\item RF-03.2. Editar reserva.
		\item RF-03.3. Eliminar reserva.
	\end{enumerate}
	\item RF-04. Generar informes.
	\begin{enumerate}
		\item RF-04.1. Informe de horarios por grado.
		\item RF-04.2. Informe de ocupación de aulas.
		\item RF-04.3. Informe de exámenes convocatorias.
	\end{enumerate}
	\item RF-05. Gestión de roles.
	\begin{enumerate}
		\item RF-05.1. Crear rol.
		\item RF-05.2. Editar rol.
		\item RF-05.3. Eliminar rol.
		\item RF-05.4. Gestionar permisos.
	\end{enumerate}
	\item RF-06. Gestión de usuarios.
	\begin{enumerate}
		\item RF-06.1. Crear usuario.
		\item RF-06.2. Editar usuario.
		\item RF-06.3. Eliminar usuario.
	\end{enumerate}
	\item RF-07. Calcular fechas de exámenes.
\end{enumerate}

\subsection{Requisitos No Funcionales}


\begin{enumerate}
	\item RNF-01. La aplicación debe ser segura.
	\item RNF-02. La aplicación debe tener un uso sencillo.
	\item RNF-03. Escalabilidad.
	\item RNF-04. Rendimiento.
	\item RNF-05. Disponibilidad.
	\item RNF-06. Interoperabilidad o compatibildad.
	\item RNF-07. Mantenimiento.
\end{enumerate}
\section{Especificación de requisitos}

\section{Especificación de casos de uso}

\subsection{Diagrama general de casos de uso}

\imagen{img/Diagrama de casos de uso.png}{Diagrama general de casos de uso}

\subsection{Definición de los casos de uso}

\begin{table}[p]
	\centering
	\begin{tabularx}{\linewidth}{ p{0.21\columnwidth} p{0.71\columnwidth} }
		\toprule
		\textbf{CU-1}    & \textbf{Inicio de sesión}\\
		\toprule
		\textbf{Versión}              & 1.0    \\
		\textbf{Autor}                & Rebeca Cuesta Estépar \\
		\textbf{Requisitos asociados} & RF-01 \\
		\textbf{Descripción}          & Un usuario no autenticado, introduce las credenciales (correo y contraseña). Si son credenciales válidas , se le permite acceder a la aplicación redirigiéndole en función de su rol, si las credenciales son inválidas se mostrará el mensaje de error y no se le permitirá el acceso. \\
		\textbf{Precondición}         & El usuario está registrado y su cuenta se encuentra activa \\
		\textbf{Acciones}             & 
		\begin{enumerate}
			\def\labelenumi{\arabic{enumi}.}
			\tightlist
			\item Introducir correo
			\item Introducir contraseña
			\item El sistema valida el formato
			\item El sistema verifica las credenciales
			\item Se crea la sesión
			\item Se redirige a la pantalla correspondiente según su rol
			
		\end{enumerate}\\
		\textbf{Postcondición}        & En caso de éxito en el inicio de sesión, el usuario tendrá una sesión activa y el usuario queda en su panel correspondiente. En caso de fallo no se crea la sesión, se incrementa el contador de intentos fallidos y se muestra el mensaje de error. \\
		\textbf{Excepciones}          & 
		\begin{itemize}
			\item\textbf{Ex1} Credenciales incorrectas \textrightarrow{} mensaje genérico y sumar el intento fallido.
			\item\textbf{Ex2} Usuario inactivo o bloqueado \textrightarrow{} informar del fallo
		\end{itemize} \\
		\textbf{Importancia}          & Alta \\
		\bottomrule
	\end{tabularx}
	\caption{CU-1 Inicio de sesión}
\end{table}

\vspace*{-2cm}
\setlength{\textfloatsep}{5pt}
\setlength{\intextsep}{5pt}
\begin{table}[p]
	\centering
	\begin{tabularx}{\linewidth}{ p{0.21\columnwidth} p{0.71\columnwidth} }
		\toprule
		\textbf{CU-2}    & \textbf{Gestionar usuarios}\\
		\toprule
		\textbf{Versión}              & 1.0    \\
		\textbf{Autor}                & Rebeca Cuesta Estépar \\
		\textbf{Requisitos asociados} &  \\
		\textbf{Descripción}          & El administrador puede crear, modificar y eliminar cuentas de usuario del sistema. Incluye búsqueda y listado de los usuarios. \\
		\textbf{Precondición}         & El administrador tiene que estar autenticado, tiene que disponer de los permisos para la gestión de usuarios y la base de datos tiene que estar disponible. \\
		\textbf{Acciones}             &
		{\small
		\begin{enumerate}
			\def\labelenumi{\arabic{enumi}.}
			\tightlist
			\item El administrador accede a la Gestión de Usuarios
			\item Visualiza el listado de usuarios (nombre, email, estado), y las opciones (Añadir, editar y eliminar)
			\item El administrador pulsa \emph{Añadir}
			\begin{enumerate}[label*=\arabic*., itemsep=0.5pt, topsep=0.5pt, leftmargin=1em]
				\item El sistema muestra un formulario con todos los campos correspondientes a los datos necesarios
				\item El adminsitrador completa estos campos
				\item El sistema valida el formato del email y que este sea único
				\item El sistema crea la cuenta con un estado y contraseña inicial (esta es aleatoria obligar a restablecer)
				\item El sistema muestra la confirmación
			\end{enumerate}
			\item El administrador selecciona al usuario y pulsa \emph{Editar}
			\begin{enumerate}[label*=\arabic*., itemsep=2pt, topsep=2pt, leftmargin=1em]
				\item El sistema muestra los datos actuales
				\item El administrador modifica los campos 
				\item El sistema valida los campos modificados
				\item El sistema guarda los cambios
			\end{enumerate}
			\item El administrador pulsa \emph{Eliminar} sobre un usuario
			\begin{enumerate}[label*=\arabic*., itemsep=2pt, topsep=2pt, leftmargin=1em]
				\item El sistema solicita confirmación y muestra cambios asociados
				\item El administrador confirma
				\item El sistema ejecuta el borrado, y elimina los datos asociados
			\end{enumerate}
		\end{enumerate}
		}
		\\
		\textbf{Postcondición}        & En caso de éxito el usuario se crea, edita o elimina según la acción realizada. En caso de fallo no se realizan los cambios y se muestra el error \\
		\textbf{Excepciones}          & 
		{\small
		\begin{itemize}
			\setlength{\itemsep}{0.1em}
			\setlength{\parsep}{0pt}
			\setlength{\topsep}{0pt}
			\tightlist
			\item \textbf{Ex1} Email duplicado \textrightarrow{} Impedir la creación o edición y mostrar el error
			\item \textbf{Ex2} Formato de datos inválidos \textrightarrow{} mostrar los errores de validación

		\end{itemize}
		} \\
		\textbf{Importancia}          & Alta \\
		\bottomrule
	\end{tabularx}
	\caption{CU-2 Gestionar usuarios}
\end{table}

\begin{table}[p]
	\centering
	\begin{tabularx}{\linewidth}{ p{0.21\columnwidth} p{0.71\columnwidth} }
		\toprule
		\textbf{CU-3}    & \textbf{Asignar roles}\\
		\toprule
		\textbf{Versión}              & 1.0    \\
		\textbf{Autor}                & Rebeca Cuesta Estépar \\
		\textbf{Requisitos asociados} &  \\
		\textbf{Descripción}          & El administrador asigna, cambio o quitar roles a usuarios para determinar los permisos de acceso. \\
		\textbf{Precondición}         & El administrador tiene que estar autenticado, que el usuario sobre el que se actúa exista, y que los roles estén disponibles. \\
		\textbf{Acciones}             &
		\begin{enumerate}
			\def\labelenumi{\arabic{enumi}.}
			\tightlist
			\item El administrador accede a \emph{Gestión de usuarios} y selecciona un usuario.
			\item El sistema muestra los datos del usuario y sus roles actuales.
			\item El administrador pulsa \emph{Añadir rol}.
			\begin{enumerate}[label*=\arabic*., itemsep=0.5pt, topsep=0.5pt, leftmargin=1em]
				\item El sistema muestra el listado de roles disponibles (no asignados).
				\item El administrador selecciona uno o varios roles y confirma.
				\item El sistema valida compatibilidad de roles.
				\item El sistema asigna el o los roles y confirma.
			\end{enumerate}
			\item El administrador pulsa \emph{Editar roles}.
			\begin{enumerate}[label*=\arabic*., itemsep=0.5pt, topsep=0.5pt, leftmargin=1em]
			\item El sistema muestra roles actuales y disponibles.
			\item El administrador desmarca o marca roles y confirma.
			\item El sistema valida y actualiza los roles y los confirma.
			\item El administrador selecciona un rol asignado y pulsa \emph{Quitar  rol}.
			\end{enumerate}
			\begin{enumerate}[label*=\arabic*., itemsep=0.5pt, topsep=0.5pt, leftmargin=1em]
			\item El sistema solicita confirmación y muestra las consecuencias
			\item El administrador confirma.
			\item El sistema quita el rol, actualiza permisos y confirma.
			\end{enumerate}
		\end{enumerate}\\
		\textbf{Postcondición}        & Se actualizan los roles del usuario, pero si falla no hay cambios y se muestra el mensaje de error \\
		\textbf{Excepciones}          & 
		{\small
			\begin{itemize}
				\tightlist
				\item \textbf{Ex1.} Usuario inactivo o inexistente \textrightarrow{} impedir acción.
				\item \textbf{Ex2.} No se puede dejar al usuario sin ningún rol \textrightarrow{} impedir quitar todos los roles
			\end{itemize} 
		}\\
		\textbf{Importancia}          & Alta \\
		\bottomrule
	\end{tabularx}
	\caption{CU-3 Asignar roles}
\end{table}


\begin{table}[p]
	\centering
	\begin{tabularx}{\linewidth}{ p{0.21\columnwidth} p{0.71\columnwidth} }
		\toprule
		\textbf{CU-4}    & \textbf{Generar informes}\\
		\toprule
		\textbf{Versión}              & 1.0    \\
		\textbf{Autor}                & Rebeca Cuesta Estépar \\
		\textbf{Requisitos asociados} &  \\
		\textbf{Descripción}          & El usuario genera los informes a partir de la información seleccionada (reservas puntuales, ocupación de aulas, horario de grados, exámenes) para poder exportar los resultados.\\
		\textbf{Precondición}         & Que sea un usuario autenticado, que tenga los privilegios necesarios para poder hacerlo, que la base de datos esté disponible y que existan datos para generar los informes.\\
		\textbf{Acciones}             &
		\begin{enumerate}
			\def\labelenumi{\arabic{enumi}.}
			\tightlist
			\item El usuario accede al módulo \emph{Generar informes}
			\item El usuario selecciona el tipo de informe que desea generar
			\item Elige los filtros que desea
			\item El sistema consulta los datos en la base de datos
			\item Se genera una vista previa del informe
			\item El sistema genera la exportación con el archivo
		\end{enumerate}\\
		\textbf{Postcondición}        & El informe se genera, y si el usuario lo selecciona, se exporta el informe. Si ocurre algún fallo se muestra el mensaje de error correspondiente \\
		\textbf{Excepciones}          & 
		\begin{itemize}
			\tightlist
			\item \textbf{Ex1:} Que no existan los datos para los filtros que ha solicitado el usuario \textrightarrow{} Mostrar el mensaje de que no hay resultados
			\item \textbf{Ex2:} Fallo en la conexión a la base de datos \textrightarrow{} mostrar el error y cancelar la generación
			\item \textbf{Ex3:} Fallo al generar el archivo de exportación \textrightarrow{} mostrar error
		\end{itemize} \\
		\textbf{Importancia}          & Alta \\
		\bottomrule
	\end{tabularx}
	\caption{CU-4 Generar informes}
\end{table}


\begin{table}[p]
	\centering
	\begin{tabularx}{\linewidth}{ p{0.21\columnwidth} p{0.71\columnwidth} }
		\toprule
		\textbf{CU-5}    & \textbf{Consultar reservas puntuales}\\
		\toprule
		\textbf{Versión}              & 1.0    \\
		\textbf{Autor}                & Rebeca Cuesta Estépar \\
		\textbf{Requisitos asociados} & RF-02 \\
		\textbf{Descripción}          & El gestor o administrador consulta las reservas que existen, y puede aplicar filtros sobre ellas y visualizar los detalles de cada una \\
		\textbf{Precondición}         & Que sea un usuario autenticado y que tenga el rol de gestor o de administrador. Que la base de datos esté disponible y que existan reservas registradas \\
		\textbf{Acciones}             &
		\begin{enumerate}
			\def\labelenumi{\arabic{enumi}.}
			\tightlist
			\item El usuario accede al módulo \emph{Consultar reservas puntuales}.
			\item El sistema muestra las reservas que existen en la base de datos
			\item El usuario puede aplicar filtros
			\item El sistema actualiza la lista con los resultados filtrados
			\item El usuario selecciona una reserva específica para ver más información
			\item El sistema muestra los detalles de la reserva
		\end{enumerate}\\
		\textbf{Postcondición}        &  Se muestra al usuario la información de las reservas con la posibilidad de realizar filtros. No se permite la modificación de datos.\\
		\textbf{Excepciones}          &  
		\begin{itemize}[itemsep=0.15em, parsep=0pt, topsep=0pt, leftmargin=*]
			\item \textbf{Ex1:} No existen reservas que cumplan los criterios → mostrar mensaje “no hay resultados”.
			\item \textbf{Ex2:} Error en la conexión a la base de datos → cancelar operación y notificar fallo.
			\item \textbf{Ex3:} Filtro inválido o mal formado → restablecer búsqueda y mostrar aviso.
		\end{itemize}\\
		\textbf{Importancia}          & Media-Alta \\
		\bottomrule
	\end{tabularx}
	\caption{CU-5 Consultar reservas puntuales}
\end{table}


\begin{table}[p]
	\centering
	\begin{tabularx}{\linewidth}{ p{0.21\columnwidth} p{0.71\columnwidth} }
		\toprule
		\textbf{CU-6}    & \textbf{Consultar ocupación de aula}\\
		\toprule
		\textbf{Versión}              & 1.0    \\
		\textbf{Autor}                & Rebeca Cuesta Estépar \\
		\textbf{Requisitos asociados} & RF-04 \\
		\textbf{Descripción}          & El administrador o el gestor consulta la ocupación del aula para una fecha específica o semanal (solo aparecen las reservas periódicas). El sistema muestra un horario con las horas que se encuentra ocupada o disponible al aula, en función del horario y de las reservas puntuales si así se requiere.\\
		\textbf{Precondición}         & Que sea un usuario autenticado con el rol de gestor o de administrador, que los horarios y las reservas se encuentren cargadas en el sistema y que la base de datos esté disponible. \\
		\textbf{Acciones}             &
		\begin{enumerate}
			\def\labelenumi{\arabic{enumi}.}
			\tightlist
			\item El usuario accede al apartado de \emph{Consultar ocupación de aula}.
			\item El sistema muestra un menú con filtros de aula, fecha o semanal.
			\item El usuario pulsa el aula y el periodo y pulsa \emph{Buscar}
			\item El sistema recoge los datos asociado a la búsqueda realizada por el usuario
			\item El sistema genera una vista con las horas ocupadas y disponibles.
		\end{enumerate}\\
		\textbf{Postcondición}        &  El sistema muestra la ocupación del aula seleccionada mostrándose claramente las horas disponibles y las ocupadas.\\
		\textbf{Excepciones}          & 
		\begin{itemize}[itemsep=0.15em, parsep=0pt, topsep=0pt, leftmargin=*]
			\tightlist
			\item \textbf{Ex1:} Que el aula no exista o no se haya seleccionado correctamente \textrightarrow mostrar mensaje de error.
			\item \textbf{Ex2:} Que no existan datos para la fecha o periodo seleccionado \textrightarrow mostrar que no hay existen reservas para ese aula y no hay datos que mostrar
			\item \textbf{Ex3:} Que la fecha introducida sea inválida \textrightarrow Mostrar mensaje de error
			\item \textbf{Ex4:} Que haya un error al recoger los datos de la consulta \textrightarrow notificar el fallo
		\end{itemize}\\ 
		\textbf{Importancia}          & Media-Alta \\
		\bottomrule
	\end{tabularx}
	\caption{CU-6 Generar informes}
\end{table}


\begin{table}[p]
	\centering
	\begin{tabularx}{\linewidth}{ p{0.21\columnwidth} p{0.71\columnwidth} }
		\toprule
		\textbf{CU-7}    & \textbf{Consultar horario de grado}\\
		\toprule
		\textbf{Versión}              & 1.0    \\
		\textbf{Autor}                & Rebeca Cuesta Estépar \\
		\textbf{Requisitos asociados} &  \\
		\textbf{Descripción}          & El gestor o administrador puede consultar el horario de grado. El sistema muestra las clases planificadas con su aula y hora.  \\
		\textbf{Precondición}         &  El usuario autenticado tiene permisos , los horarios de los grados deben de estar cargados en el sistema y la base de datos tiene que estar disponible \\
		\textbf{Acciones}             &
		\begin{enumerate}
			\def\labelenumi{\arabic{enumi}.}
			\tightlist
			\item El usuario accede al módulo \emph{Consultar horarios de grado}
			\item El sistema muestra los grados por los que filtrar.
			\item El usuario selecciona el grado del que quiere
			\item El sistema muestra filtros sobre los que aplicar 
			\item El usuario selecciona si quiere filtrar o no y pulsa \emph{Buscar}
			\item El sistema recoge los datos asociado al grado seleccionado y los filtros
		\end{enumerate}\\
		\textbf{Postcondición} & Se muestra el horario con los filtros y grado aplicados, con la informaciónn actualizada y recogiendo los datos de la base de datos        \\
		\textbf{Excepciones}          & 
		\begin{itemize}[itemsep=0.15em, parsep=0pt, topsep=0pt, leftmargin=*]
			\tightlist
			\item \textbf{Ex1:} Error de conexión a la base de datos \textrightarrow mostrar mensaje de error.
			\item \textbf{Ex2:} Que no existan horarios para los filtros seleccionados \textrightarrow mostrar mensaje de que no existen resultados
		\end{itemize} \\
		\textbf{Importancia}          & Media \\
		\bottomrule
	\end{tabularx}
	\caption{CU-7 Consultar horario de grado}
\end{table}


\begin{table}[p]
	\centering
	\begin{tabularx}{\linewidth}{ p{0.21\columnwidth} p{0.71\columnwidth} }
		\toprule
		\textbf{CU-8}    & \textbf{Cargar horario*******}\\
		\toprule
		\textbf{Versión}              & 1.0    \\
		\textbf{Autor}                & Rebeca Cuesta Estépar \\
		\textbf{Requisitos asociados} &  \\
		\textbf{Descripción}          & El administrador o gestor carga en el sistema un horario de otro año\\
		\textbf{Precondición}         & Que sea un usuario con permisos y que la base de datos se encuentre disponible para recoger los datos cargados \\
		\textbf{Acciones}             &
		\begin{enumerate}
			\def\labelenumi{\arabic{enumi}.}
			\tightlist
			\item El usuario accede a \emph{Cargar horario}
			\item El sistema muestra la opción de 
			\item Seleccionar el usuario
			\item Eliminarlo
		\end{enumerate}\\
		\textbf{Postcondición} &   Que el usuario y sus datos correspondientes se eliminen del sistema y no persistan sus datos      \\
		\textbf{Excepciones}          &  \\
		\textbf{Importancia}          & Alta \\
		\bottomrule
	\end{tabularx}
	\caption{CU-8 Cargar horario}
\end{table}


\begin{table}[p]
	\centering
	\begin{tabularx}{\linewidth}{ p{0.21\columnwidth} p{0.71\columnwidth} }
		\toprule
		\textbf{CU-9}    & \textbf{Gestión de reservas periódicas}\\
		\toprule
		\textbf{Versión}              & 1.0    \\
		\textbf{Autor}                & Rebeca Cuesta Estépar \\
		\textbf{Requisitos asociados} & RF-6.2 \\
		\textbf{Descripción}          & Un usuario con privilegios como el administrador podrá editar los datos de los usuarios  \\
		\textbf{Precondición}         & Que sea un usuario administrador y que exista el usuario y los datos del usuario que quiere editar\\
		\textbf{Acciones}             &
		\begin{enumerate}
			\def\labelenumi{\arabic{enumi}.}
			\tightlist
			\item Iniciar sesión
			\item Entrar en el listado de Usuarios
			\item Seleccionar el usuario correspondiente
			\item Editar datos
		\end{enumerate}\\
		\textbf{Postcondición} & Que los datos actualizados o modificados del usuario persistan en el sistema tras guardarlos     \\
		\textbf{Excepciones}          &  \\
		\textbf{Importancia}          & Alta \\
		\bottomrule
	\end{tabularx}
	\caption{CU-9 Gestión de reservas periódicas}
\end{table}

\begin{table}[p]
	\centering
	\begin{tabularx}{\linewidth}{ p{0.21\columnwidth} p{0.71\columnwidth} }
		\toprule
		\textbf{CU-10}    & \textbf{Asignar roles}\\
		\toprule
		\textbf{Versión}              & 1.0    \\
		\textbf{Autor}                & Rebeca Cuesta Estépar \\
		\textbf{Requisitos asociados} & RF-05 \\
		\textbf{Descripción}          & El usuario administrador podrá asignar un rol a cada usuario \\
		\textbf{Precondición}         & Que sea un usuario administrador, que exista el usuario y el rol \\
		\textbf{Acciones}             &
		\begin{enumerate}
			\def\labelenumi{\arabic{enumi}.}
			\tightlist
			\item Iniciar sesión
			\item Entrar en el listado de usuarios
			\item Seleccionar usuario
			\item Seleccionar rol
			\item Realizar la asignación
		\end{enumerate}\\
		\textbf{Postcondición} &  Que el usuario pueda realizar en el sistema lo que los permisos relacionados al rol que se le ha asignado le permitan       \\
		\textbf{Excepciones}          &  \\
		\textbf{Importancia}          & Alta \\
		\bottomrule
	\end{tabularx}
	\caption{CU-10 Asignar roles}
\end{table}


\begin{table}[p]
	\centering
	\begin{tabularx}{\linewidth}{ p{0.21\columnwidth} p{0.71\columnwidth} }
		\toprule
		\textbf{CU-11}    & \textbf{Calcular exámenes}\\
		\toprule
		\textbf{Versión}              & 1.0    \\
		\textbf{Autor}                & Rebeca Cuesta Estépar \\
		\textbf{Requisitos asociados} & RF-07 \\
		\textbf{Descripción}          & Realizar el cálculo pertinente para que la rotación de la fecha de exámenes sea la pertinente \\
		\textbf{Precondición}         & Que sea un usuario autenticado \\
		\textbf{Acciones}             &
		\begin{enumerate}
			\def\labelenumi{\arabic{enumi}.}
			\tightlist
			\item Iniciar sesión
			\item Entrar en el de Exámenes
			\item Seleccionar calcular rotación
		\end{enumerate}\\
		\textbf{Postcondición}   Que las fechas de los exámenes tras el cálculo de rotación sea la correcta     &  \\
		\textbf{Excepciones}    Los exámenes no pueden caer ni en sábado ni en domingo      &  \\
		\textbf{Importancia}          & Alta \\
		\bottomrule
	\end{tabularx}
	\caption{CU-11 Calcular exámenes}
\end{table}